\documentclass[11pt]{amsart}
\usepackage[margin=1.1in]{geometry}
\usepackage{amsmath,amssymb,amsthm,booktabs}
\usepackage[colorlinks=true,allcolors=blue]{hyperref}
\newtheorem{theorem}{Theorem}
\newtheorem{lemma}[theorem]{Lemma}
\theoremstyle{remark}
\newtheorem{remark}[theorem]{Remark}
\newcommand{\Sp}{\spadesuit}\newcommand{\He}{\heartsuit}\newcommand{\Di}{\diamondsuit}\newcommand{\Cl}{\clubsuit}
\begin{document}
\title[Poison cards in Late Binding solitaire]{There are no nine poison cards in Late Binding solitaire}
\author{Zijian Zeng}
\address{UCSI University, Kuala Lumpur, Malaysia}
\date{19 November 2026}
\begin{abstract}
In Late Binding solitaire eighteen cards are dealt in a row and the player tries to reduce them to one pile; a pile may
capture the pile next to it or three places to its left when their top cards share a suit or a rank. In Stanford's
1989 problem seminar K.~A. Ross asked for the least $n$ such that some sequence of $n$ cards makes every game that
starts with it a loss; Knuth's answer to exercise 7.2.2.2--486 of \emph{The Art of Computer Programming} records the
question, ``still unsolved,'' as whether there is a sequence of (say) nine such ``poison cards.'' There is not: every
sequence of nine cards can be completed to a winning deal. On the other hand the thirteen cards
$5\Di\,7\He\,3\Di\,6\He\,4\Sp\,A\Cl\,9\Sp\,Q\Cl\,K\Sp\,J\Cl\,2\He\,8\Di\,10\He$ are poison. So the least length
of a poison sequence is between $10$ and $13$; the seminar had believed it to be at most $17$, and perhaps provably
at most $15$.
\end{abstract}
\maketitle

\section{The game and the question}
Shuffle a deck and deal eighteen cards in a row, as eighteen piles. A \emph{capture} places a pile on top of the pile
immediately to its left, or on the pile obtained by skipping left over two other piles, provided the top cards of the
two piles have the same suit or the same rank; the combined pile takes the place of the captured one and the gap
closes. All eighteen cards are dealt before play begins (``late binding''), and the game is won if the piles can be
reduced to a single pile. Only top cards matter, so in the terms of Knuth's exercise 7.2.2.2--484 the game is the
\emph{quenching} of the sequence of the eighteen cards, where two cards are adjacent when they share a suit or a
rank~\cite[\S7.2.2.2, exercises 484--486]{knuth4b}.

The game was Problem~1 of Stanford's CS problem seminar in 1989~\cite{rossknuth}. There K.~A. Ross asked for the
minimum number $n^*$ of cards such that some initial sequence of $n^*$ cards admits no win for any continuation; the
group ``convinced itself'' that $n^*\le17$, and thought that $n^*\le15$ might be provable, but the question was not
pursued. Knuth's answer to exercise 486 calls it ``an interesting question, still unsolved'': is there a sequence of
(say) nine ``poison cards'' such that all games starting with them are lost? We call a sequence of $k\le18$ distinct
cards \emph{poison} if every deal of eighteen cards that begins with it is lost. Every extension of a poison sequence
is poison, so the poison sequences of each length $k\ge n^*$ exist and those of length $k<n^*$ do not.

\begin{theorem}\label{thm:nine}
No sequence of nine cards is poison. Hence $n^*\ge10$.
\end{theorem}

\begin{theorem}\label{thm:fourteen}
The sequence $5\Di\,7\He\,3\Di\,6\He\,4\Sp\,A\Cl\,9\Sp\,Q\Cl\,K\Sp\,J\Cl\,2\He\,8\Di\,10\He$ is poison: each of the
$39\cdot38\cdot37\cdot36\cdot35=69{,}090{,}840$ ways to deal the remaining five cards gives a lost game. Hence
$n^*\le13$.
\end{theorem}

\section{Patterns}
Permuting the four suits, or the thirteen ranks, maps the deck onto itself and preserves the relation ``same suit or
same rank,'' so it maps winning deals to winning deals. Hence whether a sequence of $k$ cards is poison depends only
on its \emph{pattern}: relabel the suits in order of first appearance and the ranks in order of first appearance. The
number of patterns of length $k$ is
\[
\begin{array}{c|cccccccccc}
k&1&2&3&4&5&6&7&8&9&10\\\hline
&1&3&15&113&1101&12657&162863&2281239&34260563&546727587
\end{array}
\]
(a pattern is a sequence of distinct pairs (suit class, rank class) in which each class first appears in increasing
order, with at most four cards per rank). So Theorem~\ref{thm:nine} amounts to $34{,}260{,}563$ claims, one winning
deal for each pattern of length nine.

\section{The computations}
\subsection*{The solver} A depth-first search over the sequences of top cards, remembering positions already known to
be lost. It also stops as soon as the ``suit--rank graph'' of the current top cards is disconnected (vertices the
suits and ranks present, an edge for each card); such a position can never be reduced to one pile, an observation
from the 1989 seminar. The solver reproduces Knuth's statements about the example deal of exercise 486 (winnable;
lost after A$\Cl\times$J$\Cl$ and after A$\Cl\times\times$10$\Cl$), and wins $46.4\%$ of $20{,}000$ random deals
(Knuth reports $4432$ wins in $10{,}000$ deals; Blake and Gent estimate $47.02\%$~\cite{gent}).

\subsection*{Theorem~\ref{thm:nine}} For each pattern of length nine the program tries random completions, half of them
chosen greedily so that each new card matches the card one or three places to its left, with an increasing node limit,
until the solver finds a win. Two complete runs with different random seeds found a winning completion for every one
of the $34{,}260{,}563$ patterns; the second run stored all the winning deals ($617$ MB).

\subsection*{Theorem~\ref{thm:fourteen}} All $69{,}090{,}840$ ordered completions were solved (in six parallel parts);
none is winnable. The thirteen cards are \emph{frozen}: no two of them at distance one or three share a suit or a
rank, so no capture is possible inside the sequence. They were found by a hill-climbing search over frozen
sequences that minimizes the number of wins among random completions; a candidate was checked exhaustively only
after it lost all of $200{,}000$ random completions. The same search found the frozen poison sequence
$7\Cl\,3\Sp\,4\Cl\,3\He\,8\Cl\,Q\Sp\,A\Cl\,9\Sp\,10\He\,2\Di\,5\Cl\,K\Sp\,J\He\,6\Di$ of fourteen cards ($1{,}771{,}560$
completions) and two poison sequences of fifteen cards,
$8\Di\,6\Sp\,9\Di\,5\He\,8\Cl\,10\Sp\,3\Cl\,K\Di\,7\He\,4\Di\,Q\He\,K\Sp\,J\He\,4\Cl\,A\Sp$ and
$3\He\,4\Sp\,9\Cl\,Q\Sp\,7\Cl\,J\He\,9\Di\,10\Sp\,4\Di\,3\Cl\,2\He\,K\Cl\,6\Di\,5\Di\,J\Sp$ ($46{,}620$ completions each).
About $8\%$ of random sequences of seventeen cards and $0.5\%$ of sixteen are poison. Two other thirteen-card
candidates that lost $20{,}000$ random completions turned out to have winning completions, found after
$612{,}807$ and $151{,}826$ completions.

\subsection*{Independent check} A separate program, written without access to the solver above, re-implemented the
game, the solver and the enumeration of patterns. It read the $34{,}260{,}563$ stored deals, found every one of them
winnable, and found that their nine-card prefixes cover all $34{,}260{,}563$ patterns, with no duplicates and none
missing; its own three counts of the patterns agree with the table above. It also confirmed that the thirteen-,
fourteen- and fifteen-card sequences above are poison, with $0$ winnable completions. Its solver agrees with a
second, brute-force solver on $20{,}000$ random deals, of which it finds $46.9\%$ winnable, and it detected five
deliberately corrupted versions of a test file.

\begin{remark}
Both solvers find about $47\%$ of random deals winnable, as do Blake and Gent~\cite{gent}; the figure of $4432$ wins in
$10{,}000$ random deals quoted in answer 486 is about five standard deviations lower.
\end{remark}

\begin{remark}
The exact value of $n^*$ is still open. Settling $n^*\ge11$ would mean finding winning completions for all
$546{,}727{,}587$ patterns of length ten, about sixteen times the work of Theorem~\ref{thm:nine}; a poison sequence of
twelve cards would need $2{,}763{,}633{,}600$ completions to be checked, about an hour on six cores; the search for
such a candidate continues.
\end{remark}

\subsection*{How this was done}
The programs were written and run by Claude (Anthropic's AI model) at the author's direction, and the independent check
by a separate AI agent that did not see the solver. No human expert has checked the results yet. Programs, logs and
data are public at \url{https://github.com/ArtificialZeng/AI-Has-Taste}.

\begin{thebibliography}{9}
\bibitem{knuth4b} D. E. Knuth, \emph{The Art of Computer Programming}, Vol.~4B: Combinatorial Algorithms, Part~2
(Addison--Wesley, 2022), \S7.2.2.2, exercises 484--486 and answers.
\bibitem{rossknuth} K. A. Ross and D. E. Knuth, \emph{A Programming and Problem Solving Seminar}, Report STAN-CS-89-1269,
Stanford University (1989), Problem~1.
\bibitem{gent} C. Blake and I. P. Gent, The winnability of Klondike Solitaire and many other patience games,
\emph{J. Artificial Intelligence Res.} \textbf{85} (2026), article 21; arXiv:1906.12314. (The Late Binding figure
$47.021\pm0.032\%$ is from I.~P. Gent's CP~2024 invited talk.)
\end{thebibliography}
\end{document}
